\subsection{Überführung in die Produktentwicklung}
\label{subsec:ueberfuehrung-produktentwicklung}

Nach der Generierung bildet das Zielverzeichnis ein eigenständiges
Produktprojekt in einem separaten Repository. Das Master-Template verantwortet
die wiederverwendbare technische Baseline und deren Ableitung; das
Produktprojekt übernimmt ab diesem Punkt Fachmodelle, Geschäftsregeln,
Benutzeroberfläche, externe Integrationen, Produktdaten und Betriebsparameter.
Hinzu kommen je nach Anforderung
Architekturentscheidungen zu Authentifizierung und Berechtigungen,
Skalierung, Monitoring, Backup, Datenschutz oder Desktop-Signing. Diese
Aufgaben werden durch die vorhandene Projektstruktur unterstützt, aber nicht
automatisch gelöst.

Master-Template und Produktprojekt besitzen anschließend getrennte
Lebenszyklen. Der untersuchte Generator kopiert ausgewählte Pfade, bietet aber
keinen dauerhaften Update- oder Synchronisationsmechanismus. Spätere
Verbesserungen des Masters gelangen daher nicht automatisch in bereits
erzeugte Projekte; sie müssen hinsichtlich Nutzen, Kompatibilität und
Produktrisiko bewertet und gezielt übernommen werden. Umgekehrt gehören
produktspezifische Änderungen nicht ohne Generalisierungs- und
Verifikationsprüfung zurück in die gemeinsame Baseline.

Schließlich sind technische Template-Abnahme und Produktabnahme zu trennen.
Die commitgebundene Abnahme belegt für die Baseline unter dokumentierten
Bedingungen unter anderem Generierung, Installation, Tests und ausgewählte
Buildpfade. Sie schließt produktive Bereitstellung, Signing sowie
produktspezifische Fach- und Sicherheitsanforderungen ausdrücklich aus
\parencite{kleiveist2026acceptance}. Die Produktabnahme muss deshalb die
konkreten Kundenanforderungen, Geschäftsregeln, Integrationen und
Betriebsbedingungen mit eigenen Tests und Nachweisen prüfen. Kapitel 800
überführt die technische Fallstudie somit in ein systematisches Auswahl- und
Generierungsverfahren, ohne daraus Kundenerfolge oder eine universelle
Eignungszusage abzuleiten.
